\begin{exercise}[Finite-action vortices and the rotational ansatz]{I-CH10-010}
Consider the positive Euclidean action of the two-dimensional Abelian Higgs
model,
\[
  S_E=\frac1{g^2}\int\dd^2\mathbf x\left[
    \frac{1}{4\mu^2}F_{ab}F_{ab}+|D_a\phi|^2
    +\mu^2V(|\phi|)\right],
  \qquad D_a=\partial_a-\ii mA_a,
\]
where $V\ge0$ has its Higgs minimum at $|\phi|=1$.  In winding sector
$n\in\mathbb Z$, use polar coordinates and the equivariant ansatz
\[
  \phi=f(r)\ee^{\ii n\theta},\qquad A_r=0,\qquad
  A_\theta^{\rm orth}=\frac{n}{mr}a(r).
\]
Derive the radial action and the boundary behavior at $r=0$ and $r=\infty$.
Show term by term that these conditions permit finite action and that the
gauge-invariant tails decay exponentially when the Higgs and vector masses are
positive.

For the symmetry statement, restrict to $|n|=1$.  Assume that the absolute
minimizer in this sector exists and is unique after its zero is fixed, up to a
gauge transformation.  Use invariance under a spatial rotation accompanied by
the compensating phase rotation of $\phi$ to prove that a representative of
the minimizer has the equivariant form above.  Explain why a symmetry-reduced
stationary point alone would give a weaker conclusion, and identify the extra
issues for higher winding or a nonstandard Higgs potential.
\end{exercise}
